\chapter{Approximation Algorithms}
\label{ch34:approx}

In Chapter~33, we mapped the boundary of computational hardness: thousands of practical problems---from vertex cover to the traveling salesperson---are NP-complete, resisting exact polynomial-time algorithms. If anybody ever finds a polynomial-time algorithm that solves one of them exactly on every input, they will have solved all of them, collected a million-dollar prize, and rewritten computer science. Most experts bet that no such algorithm exists.

So where does that leave the engineer who has to route delivery trucks tomorrow morning? You cannot wait for $2^{n}$ steps. You cannot tell your team that the problem is NP-hard and walk away: the trucks still need routes.

This chapter covers one principled response to intractability: instead of demanding the \emph{exact} optimum, we settle for an answer that is provably \emph{close} to optimal, computed in guaranteed polynomial time. We formalize the approximation ratio, design constant-factor algorithms for vertex cover and metric TSP using matchings and spanning trees, analyze logarithmic bounds for greedy set cover, and culminate in fully polynomial-time approximation schemes (FPTAS).

\begin{historybox}
In the wake of Richard Karp's 1972 proof that dozens of foundational optimization problems were NP-complete, industrial computer science confronted an urgent question: if exact polynomial solutions are intractable, can we salvage provable performance guarantees? In 1974, David S. Johnson published \emph{Approximation algorithms for combinatorial problems}, establishing the modern hierarchy of approximation: constant-factor ratios, logarithmic degradations, and polynomial-time approximation schemes (PTAS). In 1976, Nicos Christofides at Carnegie Mellon designed his celebrated $3/2$-approximation algorithm for metric TSP by computing a minimum-weight perfect matching on the odd-degree vertices of a minimum spanning tree using Edmonds' blossom algorithm. In the 1990s, the PCP Theorem (Arora, Lund, Motwani, Sudan, Szegedy, 1992) and Johan H\aa stad's 2001 breakthrough established tight inapproximability bounds: for MAX-3-SAT, approximating beyond $7/8 + \eps$ is NP-hard, meaning random coin-flipping achieves the best possible ratio unless $\mathrm{P} = \mathrm{NP}$. Prominent open questions include whether metric TSP can be approximated to $4/3$, and Subhash Khot's 2002 Unique Games Conjecture, which implies that a $2$-approximation for vertex cover cannot be improved in polynomial time.
\end{historybox}

\chapterline

\section{The one number that matters: the approximation ratio}

First we need a precise meaning for ``close to optimal.''

Fix an optimization problem -- say, find the \emph{smallest} vertex cover, or
the \emph{shortest} tour. Every input instance $I$ has some true optimal cost,
which we will write $\mathrm{OPT}(I)$, or just $C^{*}$ when the instance is
clear. We do not know $C^{*}$; if we could compute it quickly the problem
would not be NP-hard. Our algorithm, run on $I$, produces some feasible
solution of cost $C$. The whole question of this chapter is: \emph{how far
can $C$ be from $C^{*}$?}

\begin{definition}[Approximation ratio]
\label{ch34:def-ratio}
An algorithm $A$ is a \emph{$\rho$-approximation algorithm} (or
``achieves approximation ratio $\rho$'') if for \emph{every} input $I$ the
solution it returns, of cost $C$, satisfies
\[
\max\!\left(\frac{C}{C^{*}},\ \frac{C^{*}}{C}\right) \le \rho .
\]
We always have $\rho \ge 1$. For a minimization problem this reduces to
$C \le \rho\, C^{*}$; for a maximization problem it reduces to
$C \ge C^{*}/\rho$.
\end{definition}

The two-sided $\max$ in the definition is just a clever way to write one
formula that works for both kinds of problem. For \emph{minimization} you
worry about $C$ being too \emph{big}, so $C^{*}/C \le 1$ and the binding
term is $C/C^{*}$. For \emph{maximization} you worry about $C$ being too
\emph{small}, so $C/C^{*} \le 1$ and the binding term is $C^{*}/C$. Either
way $\rho$ measures the worst-case blow-up, and $\rho = 1$ would mean
``always exactly optimal.''

\begin{intuitionbox}
Think of $\rho$ as a \emph{warranty}. A $2$-approximation for vertex cover
is a promise stamped on the box: ``whatever graph you feed me, the cover I
hand back is at most twice the size of the perfect cover -- guaranteed, no
matter how adversarial your input.'' You do not get the best answer. You get
a hard ceiling on how bad the answer can be. For a problem nobody can solve
exactly in reasonable time, a hard ceiling is very valuable.
\end{intuitionbox}

One point here confuses almost everyone the first time.

\begin{trapbox}
To prove $C \le 2\,C^{*}$ you seem to need to know $C^{*}$ -- but $C^{*}$ is
exactly the thing you \emph{cannot} compute! How can you bound your distance
to a number you will never see?

The resolution drives every proof in this chapter: you never compute $C^{*}$ directly.
Instead, you find a quantity $L$ that you \emph{can} compute or reason about,
prove that $L \le C^{*}$ (so $L$ is a \emph{lower bound} on the optimum), and
separately prove that your answer satisfies $C \le 2L$. Chaining these,
\[
C \;\le\; 2L \;\le\; 2\,C^{*},
\]
and you have your guarantee \emph{without ever knowing $C^{*}$}. Almost every
proof in this chapter has exactly this shape: \emph{find a lower bound the
optimum cannot beat, then relate your output to that lower bound.}
\end{trapbox}

Figure~\ref{ch34:sandwich} draws this sandwich on a number line. It is the
skeleton of nearly every proof to come.

\begin{figure}[h]
\centering
\begin{tikzpicture}[font=\small]
  \draw[->,thick] (-0.5,0) -- (11,0) node[right] {cost};
  % tick marks
  \foreach \x/\lab/\col in {1.2/$L$/sourceblue, 4.4/$C^{*}$/targetgreen, 8.0/$C\le 2L$/trapred} {
    \draw[\col,thick] (\x,0.15) -- (\x,-0.15);
    \node[\col,below=2pt] at (\x,-0.15) {\lab};
  }
  \node[sourceblue,above=2pt] at (1.2,0.15) {\scriptsize computable};
  \node[targetgreen,above=2pt] at (4.4,0.15) {\scriptsize unknown};
  \node[trapred,above=2pt] at (8.0,0.15) {\scriptsize your output};
  % brace L <= C*
  \draw[decorate,decoration={brace,amplitude=4pt},sourceblue] (1.2,-0.75) -- (4.4,-0.75)
    node[midway,below=4pt,sourceblue]{\scriptsize $L\le C^{*}$ (lower bound)};
  % brace C <= 2L
  \draw[decorate,decoration={brace,amplitude=4pt,mirror},trapred] (1.2,0.9) -- (8.0,0.9)
    node[midway,above=4pt,trapred]{\scriptsize $C\le 2L$ (algorithm)};
\end{tikzpicture}
\caption{The lower-bound sandwich for a minimization problem. You never see
$C^{*}$ (green). You compute a lower bound $L$ (blue) that the optimum cannot
fall below, and you prove your output $C$ (red) is at most $2L$. The two
brackets together force $C \le 2L \le 2C^{*}$ without ever measuring the gap to
$C^{*}$ directly.}
\label{ch34:sandwich}
\end{figure}

\begin{focusbox}
Hold onto three things:
\begin{itemize}
\item $\rho$ is a worst-case guarantee over \emph{all} inputs, not an
average.
\item Smaller $\rho$ is better; $\rho=1$ means exact.
\item The trick to proving a ratio is to sandwich the optimum between a
computable lower bound and your output. You bound your distance to $C^{*}$
\emph{through} a proxy you can actually touch.
\end{itemize}
\end{focusbox}

A few problems are so friendly that you can drive $\rho$ as close to $1$ as
you are willing to pay for. A family of algorithms $A_{\eps}$, one for each
$\eps > 0$, where $A_{\eps}$ is a $(1+\eps)$-approximation, is called a
\emph{polynomial-time approximation scheme} (PTAS) if each $A_{\eps}$ runs in
time polynomial in the input size $n$. If the running time is \emph{also}
polynomial in $1/\eps$, it is a \emph{fully polynomial-time approximation
scheme} (FPTAS) -- the gold standard. We will build one for subset-sum at the
end of the chapter. Other problems are so hostile that, unless P $=$ NP, no
constant ratio is possible at all -- general TSP without the triangle
inequality is the classic example, as we will see.

\begin{focusbox}
\textbf{The hierarchy of approximation guarantees.} Not all NP-hard problems behave
identically under approximation. Some (like subset-sum and knapsack) admit fully
polynomial-time approximation schemes (FPTAS), achieving $(1+\eps)$ optimality for any $\eps>0$
in polynomial time. Others (like metric TSP and vertex cover) admit constant-factor approximations,
while set cover degrades logarithmically with input size. At the opposite extreme, general TSP
without the triangle inequality cannot be approximated within any polynomial ratio unless $\mathrm{P} = \mathrm{NP}$.
This material corresponds to \textbf{CLRS Chapter~35}: the vertex-cover, TSP, set-cover,
MAX-3-CNF, and subset-sum results developed here line up with sections 35.1 through 35.5
respectively. Vazirani's \emph{Approximation Algorithms} and Williamson--Shmoys'
\emph{The Design of Approximation Algorithms} provide comprehensive advanced treatments.
\end{focusbox}

\section{Vertex cover: a 2-approximation from a silly idea}
\label{ch34:vc}

A \emph{vertex cover} of an undirected graph $G=(V,E)$ is a set of vertices
$S \subseteq V$ such that every edge has at least one endpoint in $S$. Think
of vertices as guards and edges as corridors; a vertex cover is a placement
of guards so that every corridor is watched from at least one end. The
\emph{minimum vertex cover problem} asks for the smallest such $S$. It is
NP-hard, so an exact polynomial algorithm is (almost certainly) out of reach.

Here is a heuristic that looks naive but still comes with a guarantee. Repeat the following: pick \emph{any}
edge that is not yet covered, throw \emph{both} of its endpoints into your
cover, and cross off every edge touching either endpoint. Keep going until no
uncovered edges remain.

\begin{algorithm}[H]
\caption{\textsc{Approx-Vertex-Cover}$(G)$}
\KwIn{An undirected graph $G=(V,E)$.}
\KwOut{A vertex cover $S$ with $\abs{S} \le 2\,\abs{S^{*}}$.}
$S \gets \emptyset$\;
$E' \gets E$ \tcp*{the still-uncovered edges}
\While{$E' \neq \emptyset$}{
  pick any edge $(u,v) \in E'$\;
  $S \gets S \cup \{u, v\}$\;
  remove from $E'$ every edge incident to $u$ or to $v$\;
}
\Return $S$\;
\end{algorithm}

The algorithm is greedy and makes no attempt to choose \emph{which} edge to
grab, and it runs in linear
time $O(V+E)$ with the right adjacency structures. Let us watch it work.

\begin{example}[A traced run]
\label{ch34:vc-trace}
Take the graph on six vertices $\{a,b,c,d,e,f\}$ with edges
\[
\{a\text{--}b,\ b\text{--}c,\ c\text{--}d,\ c\text{--}e,\ d\text{--}e,\
d\text{--}f,\ e\text{--}f\}.
\]
Suppose the algorithm grabs edges in this order.
\begin{itemize}
\item Pick $b\text{--}c$. Add $b,c$ to $S$. This covers $a\text{--}b$,
$b\text{--}c$, $c\text{--}d$, $c\text{--}e$. Remaining: $d\text{--}e$,
$d\text{--}f$, $e\text{--}f$.
\item Pick $d\text{--}e$. Add $d,e$ to $S$. This covers $d\text{--}e$,
$d\text{--}f$, $e\text{--}f$. Remaining: none.
\end{itemize}
Output $S=\{b,c,d,e\}$, of size $4$. What is the true minimum here? It is
$\{b,d,e\}$, of size $3$, and you can check each edge against it: $a\text{--}b$
hits $b$; $b\text{--}c$ hits $b$; $c\text{--}d$ hits $d$; $c\text{--}e$ hits
$e$; $d\text{--}e$ hits both; $d\text{--}f$ hits $d$; $e\text{--}f$ hits $e$.
Every edge is covered, so $\{b,d,e\}$ is a valid cover, and a short search
confirms no cover of size $2$ exists (in fact $\{b,d,e\}$ is the \emph{only}
cover of size $3$). So the algorithm returned $4$, the optimum is $3$, and
$4 \le 2\cdot 3$ -- comfortably inside the warranty.
\end{example}

The picture below shows the key idea; with it, the proof is short.

\begin{figure}[h]
\centering
\begin{tikzpicture}[every node/.style={circle,draw,minimum size=6mm,inner sep=0pt,font=\small}]
  \node (b) at (0,0) {$b$};
  \node (c) at (2,0) {$c$};
  \node (d) at (4,0.8) {$d$};
  \node (e) at (4,-0.8) {$e$};
  \draw[very thick,structurepurple] (b)--(c);
  \draw[very thick,structurepurple] (d)--(e);
  \draw[gray] (c)--(d);
  \draw[gray] (c)--(e);
\end{tikzpicture}
\caption{The two \emph{purple} edges are the ones the algorithm picked.
They share no endpoint -- a \emph{matching}. Any cover must spend at least
one vertex on each of them, so $\abs{S^{*}} \ge 2$, while the algorithm spends
exactly $2$ vertices per picked edge.}
\label{ch34:vc-matching}
\end{figure}

\begin{theorem}[Vertex cover 2-approximation]
\label{ch34:vc-thm}
\textsc{Approx-Vertex-Cover} returns a vertex cover $S$ with
$\abs{S} \le 2\,\abs{S^{*}}$, where $S^{*}$ is a minimum vertex cover. So it
is a $2$-approximation algorithm.
\end{theorem}

\begin{proof}
First, $S$ \emph{is} a vertex cover: the loop only stops when $E'$ is empty,
and an edge leaves $E'$ exactly when one of its endpoints enters $S$. So every
edge has an endpoint in $S$. Good -- the output is feasible.

Now the ratio. Let $A$ be the \emph{set of edges} the algorithm picked in the
``pick any edge'' line (one per iteration). Here is the key observation:
\emph{no two edges in $A$ share an endpoint}. Why? When we pick $(u,v)$, we
immediately delete every edge touching $u$ or $v$ from $E'$, so no future
pick can be incident to $u$ or $v$. A set of edges with no shared endpoints is
called a \emph{matching}. So $A$ is a matching.

Two consequences fall out. First, since each picked edge contributes its two
fresh endpoints, $\abs{S} = 2\abs{A}$ exactly. Second -- and this is the lower
bound from the trap box -- \emph{any} vertex cover, including the optimal
$S^{*}$, must contain at least one endpoint of every edge in $A$. Because the
edges of $A$ are disjoint, those endpoints are all distinct, so $S^{*}$ needs
at least $\abs{A}$ vertices just to cover the matching:
\[
\abs{S^{*}} \;\ge\; \abs{A}.
\]
That is our lower bound $L = \abs{A} \le C^{*}$. Chain it:
\[
\abs{S} \;=\; 2\abs{A} \;\le\; 2\,\abs{S^{*}}. \qedhere
\]
\end{proof}

The proof follows the recipe exactly. We never computed
$\abs{S^{*}}$. We found a computable quantity ($\abs{A}$, the matching size)
that the optimum cannot undercut, and we showed our output is exactly twice it.

\begin{trapbox}
A tempting ``improvement'': instead of grabbing both endpoints of an edge,
greedily grab the single vertex of \emph{highest degree}, remove its edges,
and repeat. It sounds smarter -- and on most graphs it does better! But its
worst-case ratio is \emph{not} $2$; it degrades like $\Theta(\log n)$.
The smarter-looking rule loses the guarantee. The ``silly'' both-endpoints rule is
the one with the clean constant guarantee.
\end{trapbox}

\begin{walkbox}
The factor of $2$ comes from spending \emph{two} vertices to cover edges the
optimum might cover with \emph{one}. Can you cook up a graph where the
algorithm genuinely returns a cover of size $2\abs{S^{*}}$ -- where the bound
is \emph{tight}, rather than merely an upper bound? (Hint: a single edge already gives a
ratio of $2$ if the optimum is $1$. Now chain many disjoint edges together.)
On your walk, decide whether \emph{any} algorithm could beat ratio $2$ here.
It is a famous open question whether $2-\delta$ is possible.
\end{walkbox}

\section{The traveling salesman with a triangle inequality}

A salesman must visit $n$ cities and return home, minimizing total travel.
Formally: given a complete graph with a cost $c(u,v) \ge 0$ on each edge,
find the cheapest \emph{Hamiltonian cycle} -- a tour visiting every vertex
exactly once. This is the \emph{traveling salesman problem} (TSP), and it is
the poster child of NP-hardness.

In full generality TSP is hopeless to approximate: unless P $=$ NP, no
polynomial algorithm achieves \emph{any} constant ratio. (The reason: a good
approximation could be used to detect Hamiltonian cycles exactly, which is
NP-hard.) So we add one mild, extremely common assumption.

\begin{definition}[Triangle inequality]
\label{ch34:tri}
The costs satisfy the \emph{triangle inequality} if for all vertices
$u,v,w$,
\[
c(u,w) \;\le\; c(u,v) + c(v,w).
\]
\end{definition}

In words: a direct trip never costs more than a detour through some other
city. Any cost that comes from actual distances in a map -- Euclidean,
driving time, anything ``metric'' -- obeys this automatically. With this
assumption, we get a simple $2$-approximation built from minimum spanning
trees, which you already know well.

\begin{intuitionbox}
A \emph{minimum spanning tree} (MST) is the cheapest way to wire all the
cities together so they are connected -- no cycle, just a skeleton. Because it
connects every vertex at minimum total weight, its cost is a \emph{lower bound}
on any tour: deleting one edge from any tour leaves a spanning tree. We then
turn that skeleton into a valid tour by \emph{walking around it}, and the
triangle inequality keeps the shortcuts from increasing the cost. The MST costs
no more than the optimal tour, and converting it into a legal tour at most
doubles its weight.
\end{intuitionbox}

\begin{algorithm}[H]
\caption{\textsc{Approx-TSP-Tour}$(G,c)$}
\KwIn{Complete graph $G$ with metric costs $c$; a root vertex $r$.}
\KwOut{A Hamiltonian cycle of cost $\le 2\,C^{*}$.}
compute a minimum spanning tree $T$ of $G$ rooted at $r$ (e.g.\ Prim)\;
let $H = \langle v_1, v_2, \dots, v_n\rangle$ be the order in which a
  \emph{preorder} walk of $T$ first visits the vertices\;
\Return the cycle that visits the vertices in the order $H$ and returns to
  $v_1$\;
\end{algorithm}

The idea has three steps. (1) Build the MST $T$. (2) Imagine walking around
the outside of $T$ -- a ``full walk'' that goes down each branch and back up,
traversing every tree edge exactly \emph{twice}, costing $2\,c(T)$. (3) That
walk repeats vertices, which is illegal for a tour. So \emph{shortcut}: list
the vertices in the order the walk \emph{first} meets them (preorder) and go
directly between consecutive ones. Shortcutting can only help, thanks to the
triangle inequality.

\begin{figure}[h]
\centering
\begin{tikzpicture}[
  scale=0.88,
  >=Stealth,
  v/.style={circle, draw=sourceblue!80!black, fill=sourceblue!12, minimum size=6.5mm, inner sep=0pt, font=\small},
  tree/.style={very thick, targetgreen!80!black},
  lbl/.style={font=\sffamily\footnotesize}
]
% Panel (a): MST T
\begin{scope}[shift={(0,0)}]
  \node[v] (a) at (0,0) {$a$};
  \node[v] (b) at (-1.3,-1.2) {$b$};
  \node[v] (c) at (1.3,-1.2) {$c$};
  \node[v] (d) at (-2.0,-2.4) {$d$};
  \node[v] (e) at (-0.6,-2.4) {$e$};

  \draw[tree] (a)--(b);
  \draw[tree] (a)--(c);
  \draw[tree] (b)--(d);
  \draw[tree] (b)--(e);

  \node[lbl,align=center,below=10pt of e,xshift=-0.5cm] {\textbf{(a) MST $T$}\\$c(T) \le C^*$};
\end{scope}

% Panel (b): Full walk (doubled edges)
\begin{scope}[shift={(5.2,0)}]
  \node[v] (a2) at (0,0) {$a$};
  \node[v] (b2) at (-1.3,-1.2) {$b$};
  \node[v] (c2) at (1.3,-1.2) {$c$};
  \node[v] (d2) at (-2.0,-2.4) {$d$};
  \node[v] (e2) at (-0.6,-2.4) {$e$};

  % Doubled arrows
  \draw[->,thick,targetgreen!80!black] ([xshift=-1.5pt]a2.south west) -- ([xshift=-1.5pt]b2.north east);
  \draw[<-,thick,quotientorange] ([xshift=1.5pt]a2.south west) -- ([xshift=1.5pt]b2.north east);

  \draw[->,thick,targetgreen!80!black] ([xshift=-1.5pt]a2.south east) -- ([xshift=-1.5pt]c2.north west);
  \draw[<-,thick,quotientorange] ([xshift=1.5pt]a2.south east) -- ([xshift=1.5pt]c2.north west);

  \draw[->,thick,targetgreen!80!black] ([xshift=-1.5pt]b2.south west) -- ([xshift=-1.5pt]d2.north east);
  \draw[<-,thick,quotientorange] ([xshift=1.5pt]b2.south west) -- ([xshift=1.5pt]d2.north east);

  \draw[->,thick,targetgreen!80!black] ([xshift=-1.5pt]b2.south east) -- ([xshift=-1.5pt]e2.north west);
  \draw[<-,thick,quotientorange] ([xshift=1.5pt]b2.south east) -- ([xshift=1.5pt]e2.north west);

  \node[lbl,align=center,below=10pt of e2,xshift=-0.5cm] {\textbf{(b) Full Walk $W$}\\cost $= 2\,c(T) \le 2\,C^*$};
\end{scope}

% Panel (c): Preorder shortcutting
\begin{scope}[shift={(10.4,0)}]
  \node[v] (a3) at (0,0) {$a$};
  \node[v] (b3) at (-1.3,-1.2) {$b$};
  \node[v] (c3) at (1.3,-1.2) {$c$};
  \node[v] (d3) at (-2.0,-2.4) {$d$};
  \node[v] (e3) at (-0.6,-2.4) {$e$};

  % Preorder tour edges
  \draw[->,very thick,structurepurple] (a3) -- (b3);
  \draw[->,very thick,structurepurple] (b3) -- (d3);
  \draw[->,very thick,structurepurple,dashed] (d3) to[bend left=15] node[below,lbl,structurepurple]{skip} (e3);
  \draw[->,very thick,structurepurple,dashed] (e3) to[bend right=15] node[below=2pt,lbl,structurepurple]{skip} (c3);
  \draw[->,very thick,structurepurple] (c3) -- (a3);

  \node[lbl,align=center,below=10pt of e3,xshift=-0.5cm] {\textbf{(c) Shortcutting}\\tour cost $\le 2\,C^*$};
\end{scope}
\end{tikzpicture}
\caption{The 2-approximation for metric TSP. \textbf{(a)} The minimum spanning
tree $T$ satisfies $c(T) \le C^*$ because deleting an edge from an optimal tour
leaves a spanning tree. \textbf{(b)} Traversing every tree edge twice yields a
full walk $W = \langle a,b,d,b,e,b,a,c,a\rangle$ of cost $2\,c(T)$.
\textbf{(c)} Shortcutting to visited vertices in preorder sequence
$\langle a,b,d,e,c,a\rangle$ replaces paths with direct edges; by the triangle
inequality, direct shortcuts cannot increase the cost, giving a tour of cost
at most $2\,C^*$.}
\label{ch34:tsp-tree}
\end{figure}

\begin{theorem}[Metric TSP 2-approximation]
\label{ch34:tsp-thm}
If $c$ satisfies the triangle inequality, \textsc{Approx-TSP-Tour} returns a
tour of cost at most $2\,C^{*}$, where $C^{*}$ is the cost of an optimal tour.
\end{theorem}

\begin{proof}
\emph{Step 1 -- the MST is a lower bound.} Let $H^{*}$ be an optimal tour, of
cost $C^{*}$. Delete any single edge from $H^{*}$ and you are left with a path
through all $n$ vertices, which is a spanning tree. Since $T$ is a
\emph{minimum} spanning tree and edge costs are nonnegative,
\[
c(T) \;\le\; c(H^{*} \text{ minus one edge}) \;\le\; c(H^{*}) \;=\; C^{*}.
\]
There is our lower bound: $c(T) \le C^{*}$.

\emph{Step 2 -- the full walk costs $2\,c(T)$.} Walking down into every branch
of $T$ and back up traverses each of the $n-1$ tree edges exactly twice, for a
total of $2\,c(T)$. Call this multiset of moves $W$.

\emph{Step 3 -- shortcutting does not increase cost.} The walk $W$ visits some
vertices more than once. Our tour follows the \emph{first-visit} (preorder)
order, so wherever $W$ would revisit an already-seen vertex, our tour skips
straight to the next \emph{new} vertex. Each such skip replaces a path
$v \rightsquigarrow w$ in $W$ by the direct edge $(v,w)$; by the triangle
inequality applied repeatedly, the direct edge costs no more than the path it
replaces. Hence the tour's cost is at most the walk's cost.

Chaining the three steps:
\[
\text{tour cost} \;\le\; c(W) \;=\; 2\,c(T) \;\le\; 2\,C^{*}. \qedhere
\]
\end{proof}

The proof uses the same recipe: a computable lower bound ($c(T)$, the MST) beneath
the optimum, and an output bounded by twice it.

\begin{aestheticbox}
Metric TSP arises directly in vehicle routing, circuit
board drilling, DNA fragment assembly, warehouse pick-paths, and the order in
which a telescope slews between targets. A guaranteed-within-double tour you
can compute in $O(E\log V)$ (the cost of an MST) is the kind of thing that
runs inside logistics software shipping millions of packages a day.
And the better $1.5$-approximation below is more than academic -- variants of
it are what real solvers warm-start from.
\end{aestheticbox}

\subsection{Christofides: squeezing the ratio down to 1.5}

The factor of $2$ came entirely from \emph{doubling} the tree edges. That was
crude. The point of doubling was to make every vertex have \emph{even} degree,
because a connected graph where every vertex has even degree contains an
\emph{Eulerian circuit} -- a closed walk using each edge exactly once -- which
we can then shortcut into a tour. But doubling \emph{all} edges is overkill;
we only need to fix the vertices whose degree in $T$ is currently \emph{odd}.

Christofides' insight (1976): there is always an even number of odd-degree
vertices in $T$. Add a minimum-weight \emph{perfect matching} $M$ on just
those odd-degree vertices. Now $T \cup M$ has all even degrees, admits an
Eulerian circuit, and shortcuts to a tour. The clever part is bounding $c(M)$:
the odd vertices lie on the optimal tour in some cyclic order, which splits
into two alternating matchings; the cheaper of the two costs at most
$C^{*}/2$, and $M$ is the minimum, so $c(M) \le C^{*}/2$. Therefore
\[
\text{tour cost} \;\le\; c(T) + c(M) \;\le\; C^{*} + \tfrac{1}{2}C^{*}
\;=\; \tfrac{3}{2}\,C^{*}.
\]
Figure~\ref{ch34:christofides} contrasts the two repairs. Doubling (left)
copies \emph{every} tree edge to make all degrees even -- safe but wasteful.
Christofides (right) only adds a matching on the odd-degree vertices, a much
cheaper way to reach the same ``all degrees even'' state that admits an
Eulerian circuit.

\begin{figure}[h]
\centering
\begin{tikzpicture}[every node/.style={circle,draw,minimum size=6mm,inner sep=0pt,font=\small},scale=0.95]
  % LEFT: doubling
  \begin{scope}[xshift=0cm]
    \node (a) at (0,0) {$a$};
    \node (b) at (-1,-1.2) {$b$};
    \node (c) at (1,-1.2) {$c$};
    \node (d) at (-1,-2.6) {$d$};
    \draw[targetgreen,thick] (a)--(b);
    \draw[targetgreen,thick] (a)--(c);
    \draw[targetgreen,thick] (b)--(d);
    % doubled copies, drawn as parallel bows
    \draw[targetgreen,thick,bend left=30] (a) to (b);
    \draw[targetgreen,thick,bend left=30] (a) to (c);
    \draw[targetgreen,thick,bend left=30] (b) to (d);
    \node[draw=none] at (0,-3.4) {\scriptsize doubling: every edge twice};
  \end{scope}
  % RIGHT: christofides
  \begin{scope}[xshift=6cm]
    \node (a2) at (0,0) {$a$};
    \node (b2) at (-1,-1.2) {$b$};
    \node (c2) at (1,-1.2) {$c$};
    \node (d2) at (-1,-2.6) {$d$};
    \draw[targetgreen,thick] (a2)--(b2);
    \draw[targetgreen,thick] (a2)--(c2);
    \draw[targetgreen,thick] (b2)--(d2);
    % odd-degree vertices in this tree: a(deg2 even), b(deg2 even? a-b,b-d =>2),
    % actually leaves c,d are odd; match them
    \draw[trapred,very thick,dashed,bend right=20] (c2) to (d2);
    \node[draw=none] at (0,-3.4) {\scriptsize matching on odd vertices $\{c,d\}$};
  \end{scope}
\end{tikzpicture}
\caption{Two ways to make every degree even so an Eulerian circuit exists.
Left: double all green tree edges (cost $2\,c(T)$). Right (Christofides): the
only odd-degree vertices of this tree are the leaves $c$ and $d$; add the one
red matching edge $c\text{--}d$, costing at most $C^{*}/2$. Both reach all-even
degrees, but the matching is far cheaper, dropping the ratio from $2$ to
$3/2$.}
\label{ch34:christofides}
\end{figure}

This $1.5$-approximation stood as the best known for metric TSP from 1976
until a tiny improvement to $1.5 - 10^{-36}$ in 2020, which shows how hard
further progress has been.

\begin{complexitybox}
\textsc{Approx-TSP-Tour}: build MST in $O(E\log V)$, preorder walk in $O(V)$;
overall $O(E\log V)$, ratio $2$. Christofides: dominated by the
minimum-weight perfect matching, $O(V^{3})$, ratio $3/2$. General
(non-metric) TSP: no constant-ratio polynomial algorithm exists unless
P $=$ NP.
\end{complexitybox}

\section{Set cover: when the best you can do is logarithmic}

Now a problem where no constant ratio is known to be achievable, and where
the natural greedy algorithm's error grows -- gently -- with the input.

You are given a universe $X$ of $n$ elements and a family
$\mathcal{F} = \{S_1,\dots,S_m\}$ of subsets of $X$ whose union is all of $X$.
A \emph{set cover} is a subfamily $\mathcal{C} \subseteq \mathcal{F}$ whose
union is still all of $X$. You want the \emph{fewest} sets. Vertex cover is
the special case where each element is an edge and each set is ``the edges
touching one vertex,'' so set cover is at least as hard -- NP-hard.

The greedy algorithm is exactly what your gut suggests: repeatedly grab the
set that covers the most \emph{still-uncovered} elements.

\begin{algorithm}[H]
\caption{\textsc{Greedy-Set-Cover}$(X,\mathcal{F})$}
\KwIn{Universe $X$, family $\mathcal{F}$ covering $X$.}
\KwOut{A cover $\mathcal{C}$.}
$U \gets X$ \tcp*{the still-uncovered elements}
$\mathcal{C} \gets \emptyset$\;
\While{$U \neq \emptyset$}{
  pick $S \in \mathcal{F}$ maximizing $\abs{S \cap U}$\;
  $U \gets U \setminus S$\;
  $\mathcal{C} \gets \mathcal{C} \cup \{S\}$\;
}
\Return $\mathcal{C}$\;
\end{algorithm}

\begin{example}[A traced run]
\label{ch34:sc-trace}
Let $X=\{1,2,3,4,5,6,7,8,9,10,11,12\}$ and
\[
S_1=\{1,2,3,4,5,6\},\quad S_2=\{5,6,7,8,9\},\quad S_3=\{1,4,7,10\},
\]
\[
S_4=\{2,5,8,11\},\quad S_5=\{3,6,9,12\}.
\]
Greedy picks $S_1$ first ($6$ new elements -- the most). Uncovered now
$\{7,8,9,10,11,12\}$. Next $S_2$ covers $\{7,8,9\}$, three new; $S_3$ covers
$\{7,10\}$, two; $S_4$ covers $\{8,11\}$; $S_5$ covers $\{9,12\}$. Greedy takes
$S_2$. Uncovered $\{10,11,12\}$. Now $S_3,S_4,S_5$ each cover exactly one
remaining element, so greedy needs all three. Result:
$\{S_1,S_2,S_3,S_4,S_5\}$, size $5$.

But the \emph{optimal} cover is $\{S_3,S_4,S_5\}$ of size $3$ -- those three
sets partition all twelve elements! Greedy picked the large overlapping sets $S_1$
and $S_2$ early, leaving three disjoint singletons to cover individually. This instance, generalized, is the classic
near-worst case showing the $\ln n$ factor is real, not an artifact of loose
analysis.
\end{example}

Figure~\ref{ch34:sc-grid} lays the twelve elements out as a grid so you can
\emph{see} the trap. The optimal cover $\{S_3,S_4,S_5\}$ is the three green
\emph{rows}, which tile the grid perfectly. Greedy ignores that structure and
grabs the two fat sets $S_1,S_2$ (the shaded left-hand columns) first, after
which it is stuck mopping up the leftovers one fragment at a time.

\begin{figure}[h]
\centering
\begin{tikzpicture}[font=\scriptsize,
   cell/.style={draw,minimum size=8mm,inner sep=0pt}]
  % rows are S3,S4,S5; columns 1..4 hold their four elements in covered-order
  \matrix (m) [matrix of nodes,nodes={cell},column sep=-\pgflinewidth,row sep=-\pgflinewidth,
     ampersand replacement=\&] {
     1 \& 4 \& 7 \& 10 \\
     2 \& 5 \& 8 \& 11 \\
     3 \& 6 \& 9 \& 12 \\
  };
  % shade the first two columns to suggest the fat sets S1 (cols1-2) and S2 overlap
  \begin{scope}[on background layer]
    \fill[sourceblue!18] ($(m-1-1.north west)$) rectangle ($(m-3-2.south east)$);
  \end{scope}
  % green brackets around each optimal row
  \foreach \r/\lab in {1/$S_3$,2/$S_4$,3/$S_5$} {
    \draw[targetgreen,thick] ($(m-\r-1.north west)+(0.04,-0.04)$) rectangle ($(m-\r-4.south east)+(-0.04,0.04)$);
    \node[targetgreen,right=6pt] at (m-\r-4.east) {\lab};
  }
  \node[sourceblue] at ($(m-1-1.north)!0.5!(m-1-2.north)+(0,0.45)$) {fat sets $S_1,S_2$ live here};
\end{tikzpicture}
\caption{Set-cover trap as a grid. The optimum is the three green rows
$S_3=\{1,4,7,10\}$, $S_4=\{2,5,8,11\}$, $S_5=\{3,6,9,12\}$, which partition all
twelve elements into a clean $3$-set cover. Greedy instead seizes the fat sets
$S_1=\{1,\dots,6\}$ and $S_2=\{5,\dots,9\}$ (concentrated in the shaded
left-hand columns) because they cover the most at first, then must clean up
with all of $S_3,S_4,S_5$ anyway -- five sets where three would do.}
\label{ch34:sc-grid}
\end{figure}

\begin{theorem}[Greedy set cover ratio]
\label{ch34:sc-thm}
\textsc{Greedy-Set-Cover} returns a cover of size at most
$H(d)\cdot\abs{\mathcal{C}^{*}}$, where $d=\max_i \abs{S_i}$ is the size of the
largest set, $\mathcal{C}^{*}$ is an optimal cover, and
$H(d)=\sum_{k=1}^{d} 1/k \le \ln d + 1$ is the $d$-th harmonic number. In
particular it is an $(\ln n + 1)$-approximation.
\end{theorem}

\begin{proof}[Proof idea (the charging argument)]
This proof uses a different device: \emph{charging}. When greedy
picks a set $S$ that newly covers $t$ elements, it pays a cost of $1$ (one
set). Spread that cost evenly over the newly covered elements: charge each of
those $t$ elements a cost of $1/t$. Then the total number of sets greedy picks
equals the total charge over all elements:
$\abs{\mathcal{C}} = \sum_{x\in X} \mathrm{cost}(x)$.

Now bound the total charge by looking through the lens of the \emph{optimal}
cover. Take any optimal set $S^{*}\in\mathcal{C}^{*}$ with its elements ordered
by when greedy covered them. When greedy was about to cover the $k$-th
element of $S^{*}$, that set $S^{*}$ still had at least $\abs{S^{*}}-k+1$
uncovered elements available to greedy -- so the set greedy actually chose
covered \emph{at least} that many, meaning the per-element charge was at most
$1/(\abs{S^{*}}-k+1)$. Summing over the elements of $S^{*}$ telescopes into a
harmonic number $H(\abs{S^{*}}) \le H(d)$. Since the optimal sets cover all of
$X$, summing over $\mathcal{C}^{*}$ gives
\[
\abs{\mathcal{C}} \;=\; \sum_{x} \mathrm{cost}(x)
\;\le\; \sum_{S^{*}\in\mathcal{C}^{*}} H(d)
\;=\; H(d)\,\abs{\mathcal{C}^{*}}. \qedhere
\]
\end{proof}

\begin{focusbox}
Set cover teaches two new lessons. First, not every problem admits a constant
ratio; sometimes the true answer is $\Theta(\log n)$, and -- by a deep
result built on the PCP theorem -- that is provably the \emph{best possible}
unless P $=$ NP. Second, the \emph{charging argument} is a reusable proof
technique: distribute the algorithm's cost onto the elements, then re-collect it
through the optimum's structure.
\end{focusbox}

\section{MAX-3-CNF: how a coin flip guarantees 7/8}

The next algorithm does not look at the input at all: it simply
\emph{flips coins}.

A \emph{3-CNF formula} is an AND of \emph{clauses}, each clause an OR of
exactly three \emph{literals}, where a literal is a variable or its negation,
and -- to keep things clean -- the three literals in a clause involve three
\emph{distinct} variables. Example:
\[
(x_1 \lor \lnot x_2 \lor x_3)\ \land\
(\lnot x_1 \lor x_2 \lor x_4)\ \land\
(x_2 \lor \lnot x_3 \lor \lnot x_4).
\]
Deciding whether \emph{all} clauses can be satisfied at once is 3-SAT, the
canonical NP-complete problem. The optimization cousin, \emph{MAX-3-CNF},
asks: over all truth assignments, what is the maximum number of clauses you
can make true? This is a \emph{maximization} problem, so our ratio guarantee
will read $C \ge C^{*}/\rho$.

Here is the algorithm in its entirety: set each variable to \texttt{true} or
\texttt{false} independently with probability $1/2$ each. The algorithm never reads
the formula.

\begin{theorem}[Randomized MAX-3-CNF]
\label{ch34:max3-thm}
Given a 3-CNF formula with $m$ clauses (each over three distinct variables),
the random assignment satisfies, \emph{in expectation}, $\tfrac{7}{8}m$
clauses. Since the optimum can satisfy at most $m$, this is a randomized
$8/7$-approximation.
\end{theorem}

\begin{proof}
Fix a clause $C_j$, say $(\ell_1 \lor \ell_2 \lor \ell_3)$ on three distinct
variables. The only way the clause comes out \emph{false} is if all three
literals are false simultaneously. Each literal is false with probability
$1/2$, and because the three variables are distinct, the three events are
\emph{independent}, so
\[
\Pr[C_j \text{ unsatisfied}] = \left(\tfrac12\right)^3 = \tfrac18,
\qquad
\Pr[C_j \text{ satisfied}] = 1 - \tfrac18 = \tfrac78.
\]
Define an indicator $Y_j = 1$ if $C_j$ is satisfied, else $0$. Then
$\mathbb{E}[Y_j] = \Pr[C_j\text{ satisfied}] = 7/8$. Let $Y=\sum_j Y_j$ be the
number of satisfied clauses. By \emph{linearity of expectation} -- which needs
no independence between clauses at all --
\[
\mathbb{E}[Y] = \sum_{j=1}^{m} \mathbb{E}[Y_j] = \tfrac78 m.
\]
The optimum $C^{*}\le m$, so $\mathbb{E}[Y] = \tfrac78 m \ge \tfrac78 C^{*}$,
i.e.\ $C^{*}/\mathbb{E}[Y] \le 8/7$. \qedhere
\end{proof}

\begin{intuitionbox}
Why does blind luck do so well? Each clause is hard to falsify: you must get
\emph{all three} of its literals wrong at once, a $1/8$ event, so a random
assignment satisfies it seven times out of eight. Linearity of
expectation does the real work here -- it lets us add up the per-clause $7/8$
probabilities even though the clauses overlap in messy, dependent ways. We
never untangle the dependence; we just sum expectations.
\end{intuitionbox}

\begin{remark}[From expectation to certainty]
``In expectation'' might feel like a dodge -- what if your particular coin
flips land badly? Two fixes. (1) Run it a few times and keep the best; since
the average is $\tfrac78 m$, some run hits at least $\tfrac78 m$, and the
probability of repeatedly falling short shrinks geometrically. (2) More
elegantly, \emph{derandomize} by the \emph{method of conditional
expectations}: fix the variables one at a time, each time choosing the value
that keeps the conditional expected number of satisfied clauses at least
$\tfrac78 m$. This is always possible (one of the two choices must be at least
the average) and yields a \emph{deterministic} assignment satisfying at least
$\tfrac78 m$ clauses, no luck required.
\end{remark}

\begin{walkbox}
The $7/8$ ratio is notable: H\aa stad proved in 2001 that for MAX-3-SAT no
polynomial algorithm beats $7/8 + \delta$ unless P $=$ NP. So a coin
flip is, up to an arbitrarily small margin, the \emph{best possible} polynomial-time
algorithm. On your walk, consider that the best efficient strategy for this
NP-hard problem ignores the formula entirely. When is randomness a feature rather than a fallback?
\end{walkbox}

\section{Subset-sum: an FPTAS, the gold standard}

We close with the most refined possibility: a single dial $\eps$ you can turn
to make the answer as close to optimal as you wish -- $1\%$, $0.01\%$,
anything -- with the running time growing only \emph{polynomially} in
$1/\eps$. That is a \emph{fully polynomial-time approximation scheme}, and
subset-sum is the cleanest place to meet one.

\textbf{The problem.} Given positive integers $x_1,\dots,x_n$ and a target
$t$, find a subset whose sum is as large as possible \emph{without exceeding}
$t$. (Think: a knapsack where every item's value equals its weight, and the
sack holds $t$.) NP-hard in general.

\textbf{The exact-but-slow core.} Build the set $P_i$ of \emph{all} subset
sums achievable using the first $i$ items and not exceeding $t$. Start with
$P_0=\{0\}$; to get $P_i$, take every sum $s$ in $P_{i-1}$ and also offer
$s+x_i$, keeping only values $\le t$. The answer is the largest element of
$P_n$. Correct -- but $\abs{P_i}$ can double each step, blowing up to $2^n$
sums. Exponential.

\textbf{The fix: trimming.} The lists explode because they hold many sums that
are nearly equal. So after each step, \emph{trim}: walk the sorted list and
delete any value that is within a factor $(1+\delta)$ of a value you are
keeping. Formally, you keep a sparse ``representative'' list where consecutive
kept values differ by more than a factor $1+\delta$. Trimming loses a little
accuracy but keeps the lists short, because there are only so many powers of
$(1+\delta)$ between $1$ and $t$.

\begin{algorithm}[H]
\caption{\textsc{Approx-Subset-Sum}$(x_1,\dots,x_n,\,t,\,\eps)$}
\KwIn{Positive integers $x_i$, target $t$, accuracy $\eps\in(0,1)$.}
\KwOut{A sum $z^{*}\le t$ with $z^{*} \ge (1-\eps)\,C^{*}$.}
$L_0 \gets \langle 0 \rangle$\;
$\delta \gets \eps/(2n)$\;
\For{$i \gets 1$ \KwTo $n$}{
  $L_i \gets$ \textsc{Merge}$(L_{i-1},\ L_{i-1}+x_i)$ \tcp*{add $x_i$ to each, merge sorted}
  $L_i \gets$ \textsc{Trim}$(L_i,\ \delta)$ \tcp*{drop near-duplicates within $1+\delta$}
  remove from $L_i$ every element $> t$\;
}
\Return the largest element $z^{*}$ of $L_n$\;
\end{algorithm}

\textsc{Trim}$(L,\delta)$ scans the sorted $L$ keeping a running ``last kept''
value $\mathrm{last}$; it keeps the next element $y$ only if
$y > \mathrm{last}\cdot(1+\delta)$, then updates $\mathrm{last}\gets y$. So
every dropped element sits within a $(1+\delta)$ factor of a kept one.

\begin{example}[Trimming in action]
\label{ch34:ss-trim}
Suppose after a merge the sorted list is
\[
\langle 10,\,11,\,12,\,15,\,20,\,21,\,22,\,23,\,24,\,29\rangle,
\qquad \delta = 0.1 .
\]
Trim keeps $10$, then checks $11$: since
$10\cdot1.1=11$ and $11 \not> 11$, drop $11$; check $12$: $12 > 11$, keep it;
$12\cdot1.1=13.2$, so $15$ is kept; $15\cdot1.1=16.5$, so $20$ is kept;
$20\cdot1.1=22$, so $21$ and $22$ are dropped, $23$ kept;
$23\cdot1.1=25.3$, so $24$ is dropped, $29$ kept. Result:
$\langle 10,12,15,20,23,29\rangle$. The ten-element list became six, and every
deleted value is within $10\%$ of a survivor. Over many steps this is what
stops exponential blow-up. Figure~\ref{ch34:trim} shows the same trim on a
number line.
\end{example}

\begin{figure}[h]
\centering
\begin{tikzpicture}[font=\scriptsize,x=0.45cm]
  % number line; value v sits at x-coordinate (v-9) in the x=0.45cm frame
  \draw[->,thick] (0,0) -- (22,0);
  % kept values (green, above) and dropped (red, below)
  \foreach \v in {10,12,15,20,23,29} {
    \pgfmathsetmacro\xv{\v-9}
    \draw[targetgreen,thick] (\xv,0) -- (\xv,1);
    \node[targetgreen,above,inner sep=1pt] at (\xv,1) {$\v$};
  }
  \foreach \v in {11,21,22,24} {
    \pgfmathsetmacro\xv{\v-9}
    \draw[trapred,thick,dashed] (\xv,0) -- (\xv,-1);
    \node[trapred,below,inner sep=1pt] at (\xv,-1) {$\v$};
  }
  % show the (1+delta) reach from kept value 20 out to 22
  \draw[<->,sourceblue] (11,-2.1) -- (13,-2.1);
  \node[sourceblue,below,inner sep=1pt] at (12,-2.1) {$\times1.1$ reach};
\end{tikzpicture}
\caption{Trimming with $\delta=0.1$. Green ticks (above) are kept; red dashed
ticks (below) are dropped because each lies within a factor $1+\delta=1.1$ of
the previous kept value. For instance, keeping $20$ shadows out everything up
to $20\cdot1.1=22$, so $21$ and $22$ vanish and the next survivor is $23$. The
kept values are spread by more than a factor $1.1$ apart, which is what keeps
the list short.}
\label{ch34:trim}
\end{figure}

\begin{theorem}[Subset-sum FPTAS]
\label{ch34:ss-thm}
With $\delta=\eps/(2n)$, \textsc{Approx-Subset-Sum} returns $z^{*}$ with
$(1-\eps)\,C^{*} \le z^{*} \le C^{*}$, and runs in time polynomial in both $n$
and $1/\eps$. Hence it is an FPTAS.
\end{theorem}

\begin{proof}[Why it is accurate]
The key invariant: \emph{for every achievable sum $y\le t$ that the exact
algorithm would have in $P_i$, the trimmed list $L_i$ contains some
representative $y'$ with}
\[
\frac{y}{(1+\delta)^{i}} \;\le\; y' \;\le\; y .
\]
Each trim can shrink a value's nearest representative by at most a factor
$1+\delta$, and we trim once per item, so after $n$ items the accumulated loss
is at most $(1+\delta)^{n}$. Apply the invariant to $y = C^{*}$, the true
optimum: the returned $z^{*}$ satisfies $z^{*} \ge C^{*}/(1+\delta)^{n}$. Now
bound the denominator using $1+\delta \le e^{\delta}$:
\[
(1+\delta)^{n} \le e^{\delta n} = e^{\eps/2} \le 1 + \eps,
\]
where the last step holds for $\eps\in(0,1]$ (since $e^{u}\le 1+2u$ for
$u\in[0,1/2]$). Therefore
\[
z^{*} \;\ge\; \frac{C^{*}}{1+\eps} \;\ge\; (1-\eps)\,C^{*}.
\]
And $z^{*}\le C^{*}$ because every value in every list is a genuine subset
sum that never exceeds $t$. So the answer is sandwiched as claimed.
\end{proof}

\begin{proof}[Why it is fast]
After trimming, consecutive kept values in $L_i$ differ by more than a factor
$1+\delta$. All positive values lie in $[1,t]$, so the number of kept values
(plus the element $0$) is at most
\[
\log_{1+\delta} t + 2 \;=\; \frac{\ln t}{\ln(1+\delta)} + 2
\;\le\; \frac{2\ln t}{\delta} + 2,
\]
using $\ln(1+\delta)\ge \delta/2$ for small $\delta$. Substituting
$\delta = \eps/(2n)$ makes each list length $O\!\big(\tfrac{n\ln t}{\eps}\big)$.
Each of the $n$ iterations does work proportional to the list length, so the
total is $O\!\big(\tfrac{n^{2}\ln t}{\eps}\big)$ -- polynomial in $n$, in
$\log t$ (the \emph{number of bits} of the input), and in $1/\eps$. That
triple ``poly'' is exactly the FPTAS promise. \qedhere
\end{proof}

\begin{complexitybox}
\textsc{Approx-Subset-Sum}: each list has $O\!\big(\tfrac{n}{\eps}\ln t\big)$
entries; total time $O\!\big(\tfrac{n^{2}}{\eps}\ln t\big)$. Turning $\eps$
down from $0.1$ to $0.001$ costs a factor of $100$ in time -- expensive, but
\emph{polynomial}, never exponential. That is what ``fully'' means
in FPTAS: $1/\eps$ lives in the polynomial, not the exponent.
\end{complexitybox}

\begin{labbox}
Trace \textsc{Approx-Subset-Sum} by hand on $x=\langle 1, 4, 5\rangle$,
$t=10$, $\eps=0.4$ (so $\delta = 0.4/6 \approx 0.0667$).
\begin{enumerate}
\item $L_0=\langle 0\rangle$.
\item Item $1$: merge gives $\langle 0,1\rangle$; trim and cap at $10$ keep
both. $L_1=\langle 0,1\rangle$.
\item Item $4$: merge $\langle 0,1\rangle$ with $\langle 4,5\rangle$ gives
$\langle 0,1,4,5\rangle$; trimming changes nothing ($5/4 > 1.0667$).
$L_2=\langle 0,1,4,5\rangle$.
\item Item $5$: merge with $\langle 5,6,9,10\rangle$ gives
$\langle 0,1,4,5,6,9,10\rangle$; the two copies of $5$ collapse; trim and cap
keep the spread-out values. $L_3$ contains $10$.
\end{enumerate}
The true optimum is $1+4+5=10=t$, and the algorithm returns $10$. Now redo it
with $t=9$: the optimum is $5+4=9$, and watch whether trimming ever costs you
the exact answer. Convince yourself the returned value never drops below
$(1-\eps)\,C^{*}$.
\end{labbox}

\section{Stepping back: a field of guarantees}

We have seen five problems and five different ways to prove ``close
enough,'' but the proofs share one idea.

\begin{itemize}
\item \textbf{Vertex cover} and \textbf{metric TSP}: a \emph{constant} ratio
of $2$, each proved by sandwiching the optimum against a structure you can
compute -- a maximal matching, an MST.
\item \textbf{Set cover}: only a \emph{logarithmic} ratio $\Theta(\log n)$ is
possible, proved by a charging argument, and that logarithm is provably the
floor.
\item \textbf{MAX-3-CNF}: a \emph{coin flip} achieves a $7/8$ fraction in
expectation, via linearity of expectation, and $7/8$ is essentially optimal.
\item \textbf{Subset-sum}: a \emph{tunable} $(1+\eps)$ ratio via an FPTAS,
proved by tracking accumulated rounding error against the optimum.
\end{itemize}

Every single proof obeyed the rule from the very first trap box: \emph{find a
quantity you can compute that the optimum cannot beat, then relate your output
to that quantity}. Matching size, MST cost, harmonic charges, expected
satisfied clauses, accumulated trim error -- all proxies for the unreachable
$C^{*}$. That is the one idea to carry out of this chapter and into CLRS
Chapter 35.

\begin{aestheticbox}
Many of these ratios are not just the limits of current techniques. The PCP theorem and
its descendants prove that $\ln n$ for set cover and $7/8$ for MAX-3-SAT
cannot be improved unless P $=$ NP, while $(2-\delta)$-approximation for vertex cover is ruled out under the Unique Games Conjecture. Approximation algorithms thus map the boundary of what
efficient computation can and cannot achieve.
\end{aestheticbox}

\subsection{The trap of uncertified heuristics and non-metric spaces}

Several intuitive heuristics look sensible but have poor worst-case
guarantees:
\begin{itemize}
  \item \textbf{The nearest-neighbor fallacy in TSP.} The most intuitive greedy
        heuristic for TSP is to start at a city, always travel to the nearest
        unvisited city, and finally return home. Even in Euclidean metric spaces,
        this greedy rule can leave isolated cities until the end, forcing
        long return trips. Its approximation ratio is $\Omega(\log n)$---it gets
        strictly worse as the input grows, unlike Christofides' guaranteed $1.5$.
  \item \textbf{Non-metric costs.} Inexperienced practitioners often attempt
        to run metric approximation algorithms on general cost graphs where the
        triangle inequality does not hold. But without the triangle inequality,
        \textsc{TSP} cannot be approximated within \emph{any} polynomial-time
        ratio $\rho \ge 1$ (not even $\rho = 2^n$) unless $\mathrm{P} = \mathrm{NP}$.
        A single violated metric assumption turns a tractable approximation problem
        into an intractable one.
\end{itemize}

\subsection{Practical engineering with certified bounds}

Approximation algorithms matter in practice in several ways:
\begin{itemize}
  \item \textbf{Worst-case quality guarantees.} In
        telecommunications backbone design, circuit layout, and supply-chain
        logistics, systems cannot rely on uncertified heuristics that might
        produce erratic routes. Approximation algorithms guarantee worst-case
        performance bounds that hold on every input.
  \item \textbf{Linear programming relaxations and randomized rounding.} Work
        on approximation ratios produced two standard tools:
        relaxing discrete 0-1 integer programs into continuous linear programs
        (solvable in polynomial time via Chapter~28) and rounding
        the fractional coordinates back to integers using probabilistic coin
        flips.
  \item \textbf{Sensitivity to real-world edge costs.} Metric algorithms depend
        strictly on the triangle inequality $c(u,v) \le c(u,w) + c(w,v)$ and
        symmetry $c(u,v) = c(v,u)$. Real delivery networks have one-way streets,
        congestion delays, and turn penalties that disrupt metric symmetry,
        requiring engineering adjustments.
  \item \textbf{Inapproximability plateaus.} For many foundational problems,
        known algorithms have reached provable mathematical ceilings. Under the
        Unique Games Conjecture, no polynomial-time algorithm can achieve an
        approximation ratio better than $2$ for \textsc{Vertex-Cover}---the
        matching heuristic of Section~\ref{ch34:vc} achieves the best bound
        theoretically attainable in polynomial time.
\end{itemize}

Approximation algorithms accept bounded suboptimality to recover polynomial-time guarantees. When an algorithm must handle very large workloads, another option is to divide the work across multiple processing cores. In Chapter~35, we study parallel algorithms: dynamic multithreading, the work and span laws, and the limits of concurrent execution.

\section{Exercises}

\begin{exercise}[Tight vertex cover]
\label{ch34:ex-vc-tight}
Construct a graph on which \textsc{Approx-Vertex-Cover} returns a cover of
size \emph{exactly} $2\,\abs{S^{*}}$. (Hint: $k$ disjoint edges. What is the
minimum cover? What does the algorithm pick?) Then argue that on a star graph
$K_{1,k}$ -- one center joined to $k$ leaves -- the algorithm can return a
cover of size $2$ while the optimum is $1$, so the ratio is $2$ even here.
\end{exercise}

\begin{exercise}[Order matters?]
\label{ch34:ex-vc-order}
\textsc{Approx-Vertex-Cover} says ``pick any edge.'' Show that two different
choices of edge order can yield covers of \emph{different} sizes on the same
graph, yet both still satisfy $\abs{S}\le 2\abs{S^{*}}$. Find a small graph
exhibiting this.
\end{exercise}

\begin{exercise}[Why metric matters]
\label{ch34:ex-tsp-metric}
Give a four-city instance of TSP that \emph{violates} the triangle inequality
and on which \textsc{Approx-TSP-Tour} returns a tour far worse than twice
optimal. (You only need one edge that is wildly cheaper than the detour
implies.) Explain in one sentence why the proof of
Theorem~\ref{ch34:tsp-thm} breaks.
\end{exercise}

\begin{exercise}[Preorder cost]
\label{ch34:ex-tsp-preorder}
For the MST in Figure~\ref{ch34:tsp-tree}, write out the full walk and the
preorder tour explicitly. Assuming unit-cost tree edges and a metric, give a
concrete upper bound on the tour cost in terms of $c(T)$, and check it against
the theorem.
\end{exercise}

\begin{exercise}[Greedy set cover lower bound]
\label{ch34:ex-sc-bad}
Generalize Example~\ref{ch34:sc-trace}: design an instance on $n$ elements
where the optimum is small (say $2$ sets) but greedy is forced to take
$\Theta(\log n)$ sets. Describe the construction; you do not need to write out
every set.
\end{exercise}

\begin{exercise}[MAX-2-CNF coin flip]
\label{ch34:ex-max2}
Repeat the MAX-3-CNF analysis for clauses of exactly \emph{two} distinct
literals. What fraction of clauses does a uniformly random assignment satisfy
in expectation? State the resulting approximation ratio.
\end{exercise}

\begin{exercise}[Conditional expectations]
\label{ch34:ex-cond-exp}
On the three-clause formula in the MAX-3-CNF section, run the method of
conditional expectations by hand: decide $x_1$, then $x_2$, then $x_3,x_4$,
each time choosing the value that keeps the conditional expected number of
satisfied clauses highest. Verify your final assignment satisfies all three
clauses (so this formula is fully satisfiable).
\end{exercise}

\begin{exercise}[FPTAS error tracking]
\label{ch34:ex-fptas-error}
In \textsc{Approx-Subset-Sum} we set $\delta=\eps/(2n)$. Suppose instead a
careless implementer sets $\delta=\eps$ (no division by $n$). Show that the
accuracy guarantee can fail: bound the worst-case loss factor $(1+\eps)^{n}$
and explain why it is \emph{not} within $(1-\eps)$ of optimal for large $n$.
\end{exercise}

\begin{exercise}[Trim invariant]
\label{ch34:ex-trim}
Prove the invariant used in Theorem~\ref{ch34:ss-thm}: after $i$ rounds, for
every exact achievable sum $y\le t$ there is a kept value $y'\in L_i$ with
$y/(1+\delta)^{i}\le y'\le y$. Use induction on $i$; the inductive step
combines ``merging preserves representatives'' with ``one trim loses at most a
factor $1+\delta$.''
\end{exercise}

\begin{exercise}[Scaling instead of trimming, challenge]
\label{ch34:ex-scaling}
There is an alternate FPTAS for the $0/1$ knapsack that \emph{scales item
values} down by a common factor and then solves exactly by dynamic
programming. Look up or reconstruct the idea: if values are divided by
$K=\eps\,v_{\max}/n$ and rounded down, argue that the DP runs in time
polynomial in $n$ and $1/\eps$ and loses at most an $\eps$ fraction of value.
Compare its spirit to the trimming approach here.
\end{exercise}
